% GENERATED FILE. Edit canonical lessons, then run npm run generate:books.
% canonical-source-hash: fnv1a64:7a180341165dacb2
% canonical-lessons: JA-C251-ashita-iku, JA-C251-ikanai, JA-C251-naru, JA-R251-tomorrow-and-not

\chapter{Tomorrow, and What I Will Not Do}
\label{ch:ja-tomorrow-and-what-i-will-not-do}
\hlchaptermodality{251}

\begin{chapteropening}
\textbf{By the end of this chapter:} \emph{I can say what I will do with a time word and the non-past (ashita ikimasu), what I will not do and did not do, plainly (ikanai, ikanakatta), and what something became (naru).}

Complete JA-R251-tomorrow-and-not: say what you will and will not do, and what something became.
\end{chapteropening}

\section[\texorpdfstring{\ja{あした} \ja{いく} (ashita iku)}{あした いく (ashita iku)}]{\ja{あした} \ja{いく} (ashita iku) — the future, with a time word}

\label{lesson:JA-C251-ashita-iku}



\begin{quote}
Say \textquotedblleft{}tomorrow\textquotedblright{} and \textquotedblleft{}next week\textquotedblright{}. (\emph{ashita}, \emph{raishū}.) Then say \textquotedblleft{}I go\textquotedblright{}, politely. (\emph{ikimasu}.)
\end{quote}

\subsection*{You'll want to know: no future form}
\begin{quote}
\textbf{\ja{あした} \ja{いきます。}}
\end{quote}

\begin{quote}
\emph{ashita ikimasu.} — \textquotedblleft{}I will go tomorrow.\textquotedblright{}
\end{quote}

Japanese has no separate form for the future. \textbf{\ja{いきます}} (\emph{ikimasu}) is \textquotedblleft{}I go\textquotedblright{} and also \textquotedblleft{}I will go\textquotedblright{}; the time word says which.

\begin{quote}
\textbf{\ja{あした} \ja{やまに} \ja{いきます。}}
\end{quote}

\begin{quote}
\emph{ashita yama ni ikimasu.} — \textquotedblleft{}I will go to the mountains tomorrow.\textquotedblright{}
\end{quote}

\begin{quote}
\textbf{\ja{らいしゅう} \ja{はたらきます。}}
\end{quote}

\begin{quote}
\emph{raishū hatarakimasu.} — \textquotedblleft{}I will work next week.\textquotedblright{}
\end{quote}

Without a time word, the same verb is what you do as a rule:

\begin{quote}
\textbf{\ja{まいにち} \ja{はたらきます。}}
\end{quote}

\begin{quote}
\emph{mainichi hatarakimasu.} — \textquotedblleft{}I work every day.\textquotedblright{}
\end{quote}

\subsection*{You'll want to know: the plain future}
To a friend, the dictionary form does the same job:

\begin{quote}
\textbf{\ja{えいがを} \ja{みる。}}
\end{quote}

\begin{quote}
\emph{eiga o miru.} — \textquotedblleft{}I will see a film.\textquotedblright{}
\end{quote}

\begin{quote}
\textbf{\ja{あした} \ja{えいがを} \ja{みる。}}
\end{quote}

\begin{quote}
\emph{ashita eiga o miru.} — \textquotedblleft{}I will see a film tomorrow.\textquotedblright{}
\end{quote}

That plain form is the one you already put before \textbf{\ja{つもり}} (\emph{tsumori}), \textquotedblleft{}a plan\textquotedblright{}, and it goes before \textbf{\ja{でしょう}} (\emph{deshō}), \textquotedblleft{}probably\textquotedblright{}, too:

\begin{quote}
\textbf{\ja{らいねん} \ja{うみに} \ja{いく} \ja{つもりです。}}
\end{quote}

\begin{quote}
\emph{rainen umi ni iku tsumori desu.} — \textquotedblleft{}I plan to go to the sea next year.\textquotedblright{}
\end{quote}

\begin{quote}
\textbf{\ja{あした} \ja{あめが} \ja{ふるでしょう。}}
\end{quote}

\begin{quote}
\emph{ashita ame ga furu deshō.} — \textquotedblleft{}It will probably rain tomorrow.\textquotedblright{}
\end{quote}

\subsection*{Your turn}
Say these aloud:

\begin{itemize}
\raggedright
  \item \textquotedblleft{}It will probably rain tomorrow\textquotedblright{} — \emph{ashita ame ga furu deshō}
  \item \textquotedblleft{}It will probably snow tomorrow\textquotedblright{} — \emph{ashita yuki ga furu deshō}
  \item \textquotedblleft{}It will probably rain next week\textquotedblright{} — \emph{raishū ame ga furu deshō}
\end{itemize}

\begin{tcolorbox}[breakable,colback=teal!4,colframe=teal!35!black,title={Before you move on}]
\textquotedblleft{}I will go tomorrow\textquotedblright{}, politely? (\textbf{\ja{あした} \ja{いきます}}.) To a friend? (\textbf{\ja{あした} \ja{いく}}.) What tells the listener it is the future? (The time word, \textbf{\ja{あした}}.) \textquotedblleft{}It will probably rain\textquotedblright{}? (\textbf{\ja{あめが} \ja{ふるでしょう}}.)
\end{tcolorbox}
\section[\texorpdfstring{\ja{いかない} / \ja{たべない} (ikanai / tabenai)}{いかない / たべない (ikanai / tabenai)}]{\ja{いかない} / \ja{たべない} (ikanai / tabenai) — the plain \textquotedblleft{}not\textquotedblright{}}

\label{lesson:JA-C251-ikanai}



\begin{quote}
Say, to a friend, \textquotedblleft{}I will go tomorrow\textquotedblright{}. (\emph{ashita iku}.) Then say \textquotedblleft{}I will not go\textquotedblright{}, politely. (\emph{ikimasen}.) This lesson gives the plain \textquotedblleft{}not\textquotedblright{}.
\end{quote}

\subsection*{You'll want to know: -\ja{ない}}
A friend asks whether you are going tomorrow. The plain answer:

\begin{quote}
\textbf{\ja{いかない。}}
\end{quote}

\begin{quote}
\emph{ikanai.} — \textquotedblleft{}I won't go.\textquotedblright{}
\end{quote}

\begin{quote}
\textbf{\ja{あした} \ja{いかない。}}
\end{quote}

\begin{quote}
\emph{ashita ikanai.} — \textquotedblleft{}I won't go tomorrow.\textquotedblright{}
\end{quote}

\textbf{\ja{いかない}} (\emph{ikanai}) is \textbf{\ja{いきません}} (\emph{ikimasen}) without the politeness. A verb that drops \textbf{\ja{る}} for \textbf{\ja{た}} drops it for \textbf{\ja{ない}} too: \textbf{\ja{たべる}} (\emph{taberu}) has \textbf{\ja{たべない}} (\emph{tabenai}).

\begin{quote}
\textbf{\ja{にくは} \ja{たべない。}}
\end{quote}

\begin{quote}
\emph{niku wa tabenai.} — \textquotedblleft{}I do not eat meat.\textquotedblright{}
\end{quote}

The other verbs take the sound of their last kana in the \textbf{a} row: \textbf{\ja{いく}} has \textbf{\ja{いかない}}, \textbf{\ja{のむ}} (\emph{nomu}) has \textbf{\ja{のまない}} (\emph{nomanai}), \textbf{\ja{まつ}} (\emph{matsu}) has \textbf{\ja{またない}} (\emph{matanai}). A verb in \textbf{\ja{う}} takes \textbf{\ja{わ}}: \textbf{\ja{かう}} (\emph{kau}) has \textbf{\ja{かわない}} (\emph{kawanai}).

\begin{quote}
\textbf{\ja{コーヒー}\ja{は} \ja{のまない。}}
\end{quote}

\begin{quote}
\emph{kōhī wa nomanai.} — \textquotedblleft{}I do not drink coffee.\textquotedblright{}
\end{quote}

Three are short: \textbf{\ja{する}} has \textbf{\ja{しない}} (\emph{shinai}), \textbf{\ja{くる}} has \textbf{\ja{こない}} (\emph{konai}), and \textbf{\ja{ある}} (\emph{aru}), \textquotedblleft{}there is\textquotedblright{}, has plain \textbf{\ja{ない}} (\emph{nai}).

\begin{quote}
\textbf{\ja{じかんが} \ja{ない。}}
\end{quote}

\begin{quote}
\emph{jikan ga nai.} — \textquotedblleft{}I have no time.\textquotedblright{}
\end{quote}

\subsection*{You'll want to know: -\ja{なかった}}
For \textquotedblleft{}did not\textquotedblright{}, the last \textbf{\ja{い}} of \textbf{\ja{ない}} becomes \textbf{\ja{かった}}: \textbf{\ja{いかなかった}} (\emph{ikanakatta}), \textbf{\ja{たべなかった}} (\emph{tabenakatta}). It is \textbf{\ja{ませんでした}} without the politeness.

\begin{quote}
\textbf{\ja{きのう} \ja{いかなかった。}}
\end{quote}

\begin{quote}
\emph{kinō ikanakatta.} — \textquotedblleft{}I did not go yesterday.\textquotedblright{}
\end{quote}

\begin{quote}
\textbf{\ja{あさごはんを} \ja{たべなかった。}}
\end{quote}

\begin{quote}
\emph{asagohan o tabenakatta.} — \textquotedblleft{}I did not eat breakfast.\textquotedblright{}
\end{quote}

Like the plain past, the plain \textquotedblleft{}not\textquotedblright{} has a second job that later lessons give it: it goes in front of the word that keeps a sentence polite.

\subsection*{Your turn}
Say these aloud:

\begin{itemize}
\raggedright
  \item \textquotedblleft{}I do not drink coffee\textquotedblright{} — \emph{kōhī wa nomanai}
  \item \textquotedblleft{}I do not drink tea\textquotedblright{} — \emph{ocha wa nomanai}
  \item \textquotedblleft{}I do not drink water\textquotedblright{} — \emph{mizu wa nomanai}
\end{itemize}

\begin{tcolorbox}[breakable,colback=teal!4,colframe=teal!35!black,title={Before you move on}]
\textquotedblleft{}I won't go\textquotedblright{}? (\textbf{\ja{いかない}}, \emph{ikanai}.) \textquotedblleft{}I won't buy it\textquotedblright{}? (\textbf{\ja{かわない}}, \emph{kawanai}.) \textquotedblleft{}I won't do it\textquotedblright{}? (\textbf{\ja{しない}}, \emph{shinai}.) \textquotedblleft{}He won't come\textquotedblright{}? (\textbf{\ja{こない}}, \emph{konai}.) \textquotedblleft{}I did not go\textquotedblright{}? (\textbf{\ja{いかなかった}}, \emph{ikanakatta}.)
\end{tcolorbox}
\section[\texorpdfstring{\ja{なる} (naru)}{なる (naru)}]{\ja{なる} (naru) — to become; to get}

\label{lesson:JA-C251-naru}



\begin{quote}
Before the new word: say \textquotedblleft{}I will not go\textquotedblright{}, to a friend. (\emph{ikanai}.) Then say \textquotedblleft{}cold\textquotedblright{} and \textquotedblleft{}a teacher\textquotedblright{}. (\emph{samui}, \emph{sensei}.)
\end{quote}

\subsection*{You'll want to know: \ja{なる}}
\textbf{\ja{なる}} — \emph{naru} — \textquotedblleft{}to become\textquotedblright{}, \textquotedblleft{}to get\textquotedblright{}. Polite \textbf{\ja{なります}}, \emph{narimasu}; a godan verb. Past (polite): \textbf{\ja{あには} \ja{せんせいに} \ja{なりました}}, \emph{ani wa sensei ni narimashita}.

After a noun it takes \textbf{\ja{に}}:

\begin{quote}
\textbf{\ja{あには} \ja{せんせいに} \ja{なりました。}}
\end{quote}

\begin{quote}
\emph{ani wa sensei ni narimashita.} — \textquotedblleft{}My elder brother became a teacher.\textquotedblright{}
\end{quote}

After an i-adjective it takes the \textbf{-\ja{く}} form you know from \textbf{\ja{はやく}}: \textbf{\ja{さむい}} becomes \textbf{\ja{さむく}}.

\begin{quote}
\textbf{\ja{さむく} \ja{なりました。}}
\end{quote}

\begin{quote}
\emph{samuku narimashita.} — \textquotedblleft{}It got cold.\textquotedblright{}
\end{quote}

\begin{quote}
\textbf{\ja{そらが} \ja{くらく} \ja{なりました。}}
\end{quote}

\begin{quote}
\emph{sora ga kuraku narimashita.} — \textquotedblleft{}The sky got dark.\textquotedblright{}
\end{quote}

\begin{cousinweb}
\textbf{\ja{なる}} (\emph{naru}) is a native Japanese word (wago), inherited rather than borrowed from Chinese. JMdict marks it as usually written in kana, as here, and gives \textquotedblleft{}to become\textquotedblright{} and \textquotedblleft{}to get\textquotedblright{} as its first senses (JMdict).
\end{cousinweb}

\subsection*{Your turn}
Say these aloud:

\begin{itemize}
\raggedright
  \item \textquotedblleft{}It got cold\textquotedblright{} — \emph{samuku narimashita}
  \item \textquotedblleft{}It got hot\textquotedblright{} — \emph{atsuku narimashita}
  \item \textquotedblleft{}It got dark\textquotedblright{} — \emph{kuraku narimashita}
\end{itemize}

\begin{tcolorbox}[breakable,colback=teal!4,colframe=teal!35!black,title={Before you move on}]
\textquotedblleft{}To become\textquotedblright{}? (\textbf{\ja{なる}}, \emph{naru}.) After a noun, which particle? (\textbf{\ja{に}}: \textbf{\ja{せんせいに} \ja{なりました}}.) After \textbf{\ja{さむい}}? (\textbf{\ja{さむく} \ja{なりました}}.)
\end{tcolorbox}
\section[(review)]{Review — tomorrow, and what I will not do}

\label{lesson:JA-R251-tomorrow-and-not}



\begin{quote}
Say \textquotedblleft{}I will go tomorrow\textquotedblright{}, politely, then \textquotedblleft{}I will not go\textquotedblright{}, to a friend. (\emph{ashita ikimasu}, \emph{ikanai}.)
\end{quote}

\subsection*{Your turn}
Say these aloud:

\begin{itemize}
\raggedright
  \item I will work next week.
  \item I will see a film tomorrow. — to a friend
  \item It will probably snow tomorrow.
  \item I do not eat meat, and I do not drink coffee. — to a friend
  \item My younger brother became a teacher.
\end{itemize}

\begin{tcolorbox}[breakable,colback=teal!4,colframe=teal!35!black,title={Before you move on}]
Is there a future form? (No: the non-past and a time word.) \textquotedblleft{}I won't come\textquotedblright{}, plainly? (\textbf{\ja{こない}}.) \textquotedblleft{}There is no time\textquotedblright{}? (\textbf{\ja{じかんが} \ja{ない}}.) \textquotedblleft{}I did not eat\textquotedblright{}? (\textbf{\ja{たべなかった}}.) \textquotedblleft{}It got hot\textquotedblright{}? (\textbf{\ja{あつく} \ja{なりました}}.)
\end{tcolorbox}
